\section{Conclusion}
\label{sec:conclusion}

We have demonstrated that scaling the security of \textsc{Kyber} is straightforward up to $k=7$ and 
requires generalizing the noise distribution beyond that. If future advances in cryptanalysis substantially reduce the security of 
current ML-KEM parameters, scaling the module rank provides a plausible path to restoring security margin while preserving the efficiency 
advantages of NTT arithmetic.

A hypothetical break of the module structure would still call for an unstructured scheme like FrodoKEM, so the two hedges remain complementary 
rather than substitutes. Notably, these two approaches could also be combined in a hybrid construction, deriving the final key from both a high-margin 
module KEM such as \textsc{Kaiburr} and an unstructured LWE KEM such as FrodoKEM. Such a construction could combine security-margin and structural-diversity 
hedges, at the cost of increased communication and computation.

Our noise family $f(t)$ is not specific to ML-KEM, and it is an open question whether the same construction could extend the security margin of other 
module-LWE schemes, such as ML-DSA. Finally, our ACC results show that rejection sampling becomes the bottleneck as $k$ grows; this leaves dedicated 
hardware support for rejection sampling as a direction for future work.